% Bagian: sets-functions-relations
% Bab: size-of-sets
% Subbagian: non-enumerability-alt

\documentclass[../../../include/open-logic-section]{subfiles}

\begin{document}

\olfileid[id]{sfr}{siz}{nen-alt}

\olsection{Himpunan \printtoken{S}{nonenumerable}}

\begin{editorial}
  Bagian ini membuktikan ketakterhitungan $\Bin^\omega$ dan $\Pow{\Nat}$
  dengan menggunakan definisi dalam \olref[enm-alt]{sec}, yakni dengan
  mensyaratkan bijeksi dengan~$\Nat$, bukan surjeksi dari $\PosInt$.
\end{editorial}

\begin{explain}
Himpunan $\Nat$ dari bilangan asli bersifat tak hingga. Himpunan itu juga
secara trivial !!{enumerable}. Namun, fakta yang luar biasa ialah bahwa
terdapat himpunan \emph{!!{nonenumerable}}, yakni himpunan yang tidak
!!{enumerable} (lihat \olref[sfr][siz][enm-alt]{defn:enumerable}).

Hal ini mungkin mengejutkan. Bagaimanapun, mengatakan bahwa $A$ adalah
!!{nonenumerable} berarti mengatakan bahwa \emph{tidak ada} !!{bijection} $f
\colon \Nat \to A$; yakni, tidak ada fungsi yang memetakan !!{element}s
$\Nat$ yang tak terhingga banyaknya ke~$A$ dan mencakup seluruh~$A$. Jadi,
jika $A$ !!{nonenumerable}, terdapat ``lebih banyak'' !!{element}s dari~$A$
daripada bilangan asli.

Untuk membuktikan bahwa suatu himpunan !!{nonenumerable}, Anda harus
menunjukkan bahwa tidak mungkin ada !!{bijection} yang sesuai. Cara terbaik
untuk melakukannya ialah menunjukkan bahwa setiap upaya mengenumerasi
!!{element}s~$A$ pasti melewatkan setidaknya satu !!{element}; ini menunjukkan
bahwa tidak ada fungsi $f\colon \Nat \to A$ yang !!{surjective}. Strategi umum
untuk membuktikan hal ini ialah menggunakan \emph{metode diagonalisasi} Cantor.
Jika diberikan daftar !!{element}s dari $A$, misalnya $x_1$, $x_2$, \dots,
kita membangun satu !!{element} lain dari~$A$ yang, karena cara pembangunannya,
tidak mungkin berada dalam daftar tersebut.

Namun, semua ini paling baik dipahami melalui contoh. Jadi, contoh pertama
kita adalah himpunan~$\Bin^\omega$ dari semua untai tak hingga yang terdiri
atas $0$ dan $1$. (`$\Bin$' berarti biner, dan kita cukup memandangnya sebagai
himpunan dengan dua anggota $\{0,1\}$.)
\oliflabeldef{sfr:card-arithmetic:card-opps:sec}{\footnote{Secara lebih tepat,
kita harus menetapkan bahwa $\Bin^\omega$ adalah himpunan semua barisan-$\omega$
dari $0$ dan $1$, yakni himpunan $\funfromto{\omega}{\{0,1\}}$. Namun,
maknanya baru akan menjadi jelas dalam
\olref[sfr][card-arithmetic][card-opps]{sec}.}{} Rumusan yang sedikit longgar
ini seharusnya tidak menimbulkan kebingungan untuk saat ini.}
\end{explain}

\begin{thm}
\ollabel{thm:nonenum-bin-omega}
$\Bin^\omega$~adalah !!{nonenumerable}.
\end{thm}

\begin{proof}
Untuk memperoleh kontradiksi, andaikan terdapat enumerasi seluruh
$\Bin^\omega$. Jadi, kita mempunyai daftar $s_{0}$, $s_{1}$, $s_{2}$,
\dots{} dengan setiap $s_n$
berupa untai tak hingga dari $0$ dan~$1$. Misalkan $s_n(m)$ adalah digit dalam
daftar ini dengan indeks untai $n$ dan indeks digit $m$. Dengan demikian, kita dapat memandang
daftar kita sebagai suatu larik, dengan $s_n(m)$ ditempatkan pada baris ke-$n$
dan kolom ke-$m$:
\[
\begin{array}{c|c|c|c|c|c}
& 0 & 1 & 2 & 3 & \dots \\\hline
0 & \mathbf{s_{0}(0)} & s_{0}(1) & s_{0}(2) & s_0(3) & \dots \\\hline
1 & s_{1}(0)& \mathbf{s_{1}(1)} & s_1(2) & s_1(3) & \dots \\\hline
2 & s_{2}(0)& s_{2}(1) & \mathbf{s_2(2)} & s_2(3) & \dots \\\hline
3 & s_{3}(0)& s_{3}(1) & s_3(2) & \mathbf{s_3(3)} & \dots \\\hline
\vdots & \vdots & \vdots & \vdots & \vdots & \mathbf{\ddots}
\end{array}
\]
Sekarang kita akan membangun suatu untai tak hingga $d$ dari $0$ dan $1$ yang
tidak berada dalam daftar ini. Kita melakukannya dengan menentukan setiap
entri untai tersebut, yakni kita menentukan $d(n)$ untuk semua $n \in \Nat$.
Secara intuitif, kita membaca diagonal larik di atas dari atas ke bawah
(karena itulah disebut ``metode diagonalisasi''), lalu mengubah setiap $1$ menjadi
$0$ dan setiap $0$ menjadi~$1$. Secara lebih abstrak, kita mendefinisikan
$d(n)$ sebagai $0$ atau $1$ menurut apakah !!{element} ke-$n$ pada diagonal,
$s_n(n)$, bernilai $1$ atau $0$, yaitu:
\[
d(n) =
\begin{cases}
1 & \text{jika $s_{n}(n) = 0$}\\
0 & \text{jika $s_{n}(n) = 1$}
\end{cases}
\]
Jelas bahwa $d \in \Bin^\omega$, karena ia adalah untai tak hingga dari $0$
dan $1$. Namun, kita membangun $d$ sedemikian sehingga $d(n) \neq s_n(n)$
untuk setiap $n \in \Nat$. Artinya, $d$ berbeda dari $s_n$ pada entri ke-$n$.
Jadi, $d \neq s_n$ untuk setiap $n\in \Nat$. Karena itu, $d$ tidak mungkin
berada dalam daftar $s_0$, $s_1$, $s_2$, \dots

Kita telah menunjukkan bahwa sebarang enumerasi dari suatu himpunan bagian
$\Bin^\omega$ akan melewatkan suatu !!{element} dari $\Bin^\omega$. Jadi,
tidak terdapat enumerasi dari himpunan $\Bin^\omega$; yakni, $\Bin^\omega$
adalah !!{nonenumerable}.
\end{proof}

\begin{explain}
Metode pembuktian ini disebut ``diagonalisasi'' karena menggunakan diagonal
larik untuk mendefinisikan~$d$. Namun, diagonalisasi tidak harus melibatkan
sebuah larik. Kita dapat menunjukkan bahwa suatu himpunan !!{nonenumerable}
dengan menggunakan gagasan serupa sekalipun tidak ada larik dan tidak ada
diagonal yang sesungguhnya. Hasil berikut menggambarkan caranya.
\end{explain}

\begin{thm}
\ollabel{thm:nonenum-pownat}
$\Pow{\Nat}$ tidak !!{enumerable}.
\end{thm}

\begin{proof}
Kita melanjutkan dengan cara yang sama, yaitu dengan menunjukkan bahwa setiap
daftar himpunan bagian dari~$\Nat$ melewatkan suatu himpunan bagian dari
$\Nat$. Jadi, andaikan kita mempunyai suatu daftar $N_0, N_1, N_2, \ldots$
dari himpunan bagian $\Nat$. Kita mendefinisikan suatu himpunan $D$ sebagai
berikut: $n \in D$ jika dan hanya jika $n \notin N_{n}$:
\[
D = \Setabs{n \in \Nat}{n \notin N_n}
\]
Jelas bahwa $D\subseteq \Nat$. Namun, $D$ tidak mungkin berada dalam daftar.
Menurut pembangunannya, $n \in D$ jika dan hanya jika $n\notin N_n$, sehingga
$D \neq N_n$ untuk setiap $n \in \Nat$.
\end{proof}

\begin{explain}
Bukti sebelumnya tidak menyebut diagonal. Meskipun demikian, Anda dapat
memandangnya sebagai bukti yang melibatkan diagonal dengan membayangkannya
sebagai berikut: Bayangkan himpunan $N_0$, $N_1$, \dots, ditulis dalam suatu
larik, dengan $N_n$ ditulis pada baris ke-$n$ dengan menuliskan $m$ pada kolom
ke-$m$ jika dan hanya jika $m \in N_n$. Sebagai contoh, misalkan empat
himpunan pertama dalam daftar itu ialah $\{0,1,2,\dots\}$, $\{1, 3, 5,
\dots\}$, $\{0,1,4\}$, dan $\{2,3,4,\dots\}$; larik kita akan dimulai sebagai
berikut:
\[
\begin{array}{r@{}rrrrrrr}
  N_0 = \{ & \mathbf{0}, & 1, & 2, & & & & \dots\}\\
  N_1 = \{ &  & \mathbf{1}, &  & 3, &  & 5, & \dots\}\\
  N_2 = \{ & 0, & 1, &  &  & 4\phantom{,} &  & \}\\
  N_3 = \{ &  &  & 2, & \mathbf{3}, & 4, & & \dots\}\\
  &\vdots & & & & & & \ddots\phantom{\}}
  \end{array}
\]
Kemudian $D$ adalah himpunan yang diperoleh dengan menelusuri diagonal dari
atas ke bawah, dengan memasukkan $n \in D$ jika dan hanya jika $n$ \emph{tidak}
berada pada diagonal. Jadi, dalam kasus di atas, kita akan mengabaikan $0$ dan
$1$, memasukkan~$2$, mengabaikan~$3$, dan seterusnya.
\end{explain}

\begin{prob}
Tunjukkan dengan argumen diagonalisasi eksplisit bahwa himpunan semua fungsi
$f \colon \Nat \to \Nat$ adalah !!{nonenumerable}. Artinya, tunjukkan bahwa
jika $f_1$, $f_2$, \dots, adalah suatu daftar fungsi dan setiap $f_i\colon
\Nat \to \Nat$, maka terdapat suatu $g \colon \Nat \to \Nat$ yang tidak
berada dalam daftar ini.
\end{prob}

\end{document}
